%%%%%%%%%%%%%%%%%%%%%%%%%%%%%%%%%%%%%%%%%%%%%%%%%%%%%%%%%%%%%%%%%%%%%%%%%%%%%%%
%% tikz-library.tex -- shared diagram macros for CS F301
%%
%% %%%%%%%%%%%%%%%%%%%%%%%%%%%%%%%%%%%%%%%%%%%%%%%%%%%%%%%%%%%%%%%%%%%%%%%%%%%%%%%
%% tikz-library.tex -- shared diagram macros for CS F301
%%
%% %%%%%%%%%%%%%%%%%%%%%%%%%%%%%%%%%%%%%%%%%%%%%%%%%%%%%%%%%%%%%%%%%%%%%%%%%%%%%%%
%% tikz-library.tex -- shared diagram macros for CS F301
%%
%% %%%%%%%%%%%%%%%%%%%%%%%%%%%%%%%%%%%%%%%%%%%%%%%%%%%%%%%%%%%%%%%%%%%%%%%%%%%%%%%
%% tikz-library.tex -- shared diagram macros for CS F301
%%
%% \input{common/tikz-library} once, from the preamble of slides or notes.
%%
%% Every macro here produces a figure a student can redraw by hand on an
%% answer sheet in under two minutes. If a new macro cannot meet that test,
%% it does not belong in this file.
%%%%%%%%%%%%%%%%%%%%%%%%%%%%%%%%%%%%%%%%%%%%%%%%%%%%%%%%%%%%%%%%%%%%%%%%%%%%%%%

\tikzset{
  csfbox/.style   = {draw=bitsblue, thick, fill=white,
                     minimum height=6mm, minimum width=22mm, font=\small\ttfamily},
  csfslot/.style  = {draw=bitsblue!70, fill=white,
                     minimum height=5.5mm, minimum width=16mm, font=\scriptsize\ttfamily},
  csfhead/.style  = {draw=bitsblue, fill=bitsblue!12, thick,
                     minimum height=6mm, minimum width=22mm,
                     font=\small\sffamily\bfseries},
  csfdead/.style  = {draw=errred, fill=errred!8, thick,
                     minimum height=6mm, minimum width=22mm, font=\small\ttfamily},
  csflive/.style  = {draw=okgreen, fill=okgreen!8, thick,
                     minimum height=6mm, minimum width=22mm, font=\small\ttfamily},
  csflabel/.style = {font=\scriptsize\sffamily, color=mutedgrey},
  csfarrow/.style = {-{Stealth[length=2mm]}, bitsblue, thick},
  csfdash/.style  = {-{Stealth[length=2mm]}, mutedgrey, dashed, thin},
  csfnode/.style  = {draw=bitsblue, circle, thick, fill=white,
                     inner sep=1.6pt, font=\small},
  csfleaf/.style  = {draw=none, font=\small\ttfamily, color=black},
}

%%==========================================================================
%% 1. ACTIVATION RECORDS AND THE CALL STACK   (M09, M10)
%%==========================================================================
%% \arframe{<x>}{<y>}{<name>}{<comma-list of slot labels, bottom to top>}
%% Draws one activation record with a titled header and stacked slots.
\newcounter{csf@slotcount}
\newcommand{\arframe}[4]{%
  \setcounter{csf@slotcount}{0}%
  \foreach \slot in {#4} {\stepcounter{csf@slotcount}}%
  \begin{scope}[shift={(#1,#2)}]
    \foreach \slot [count=\i from 0] in {#4} {
      \node[csfslot, anchor=south west] at (0, 0.55*\i) {\slot};
    }
    \node[csfhead, anchor=south west, minimum width=16mm]
      at (0, {0.55*\value{csf@slotcount}}) {\scriptsize #3};
  \end{scope}}

%% \stacklabel{<x>}{<y>}{<text>} -- annotate a frame from the side
\newcommand{\stacklabel}[3]{\node[csflabel, anchor=west] at (#1,#2) {#3};}

%% \stackgrows{<x>}{<ytop>}{<ybot>} -- the downward growth arrow
\newcommand{\stackgrows}[3]{%
  \draw[csfdash] (#1,#2) -- (#1,#3)
    node[midway, csflabel, rotate=90, anchor=south, yshift=2pt] {stack grows};}

%%==========================================================================
%% 2. HEAP BLOCKS AND POINTERS   (M07)
%%==========================================================================
%% \heapblock{<name>}{<x>}{<y>}{<content>}{<state: live|dead|plain>}
\newcommand{\heapblock}[5]{%
  \ifthenelse{\equal{#5}{dead}}%
    {\node[csfdead, anchor=south west] (#1) at (#2,#3) {#4};}%
    {\ifthenelse{\equal{#5}{live}}%
      {\node[csflive, anchor=south west] (#1) at (#2,#3) {#4};}%
      {\node[csfbox, anchor=south west] (#1) at (#2,#3) {#4};}}}

%% \pointsto{<from>}{<to>} and \dangling{<from>}
\newcommand{\pointsto}[2]{\draw[csfarrow] (#1) -- (#2);}
\newcommand{\dangling}[1]{%
  \draw[-{Stealth[length=2mm]}, errred, thick]
    (#1) -- ++(1.2,0) node[anchor=west, csflabel, color=errred] {dangling};}
\newcommand{\unreachable}[1]{%
  \node[csflabel, color=errred, anchor=west, xshift=3mm] at (#1.east)
    {unreachable};}

%%==========================================================================
%% 3. ARRAY LAYOUT AND ADDRESS COMPUTATION   (M05)
%%==========================================================================
%% \memrow{<x>}{<y>}{<n cells>}{<first index>}{<colour>}
%% A run of contiguous memory cells with index labels beneath.
\newcommand{\memrow}[5]{%
  \begin{scope}[shift={(#1,#2)}]
    \foreach \i in {0,...,\numexpr#3-1\relax} {
      \node[draw=#5, fill=#5!10, minimum size=7mm, font=\scriptsize\ttfamily,
            anchor=south west] at (0.7*\i, 0) {};
      \node[csflabel, anchor=north] at (0.7*\i+0.35, -0.05)
        {\pgfmathparse{int(\i+#4)}\pgfmathresult};
    }
  \end{scope}}

%% \arraycell{<x>}{<y>}{<label>}{<colour>} -- one highlighted cell
\newcommand{\arraycell}[4]{%
  \node[draw=#4, fill=#4!25, thick, minimum size=7mm,
        font=\scriptsize\ttfamily, anchor=south west] at (#1,#2) {#3};}

%% \rowmajorgrid{<rows>}{<cols>} -- grid annotated with linear offsets
\newcommand{\rowmajorgrid}[2]{%
  \foreach \r in {0,...,\numexpr#1-1\relax} {
    \foreach \c in {0,...,\numexpr#2-1\relax} {
      \pgfmathtruncatemacro{\off}{\r*#2+\c}
      \node[draw=bitsblue!70, fill=bitsblue!6, minimum size=8mm,
            font=\scriptsize\ttfamily, anchor=north west]
        at (0.8*\c, -0.8*\r) {\off};
    }}}

%% \colmajorgrid{<rows>}{<cols>}
\newcommand{\colmajorgrid}[2]{%
  \foreach \r in {0,...,\numexpr#1-1\relax} {
    \foreach \c in {0,...,\numexpr#2-1\relax} {
      \pgfmathtruncatemacro{\off}{\c*#1+\r}
      \node[draw=bitsgold!80!black, fill=bitsgold!10, minimum size=8mm,
            font=\scriptsize\ttfamily, anchor=north west]
        at (0.8*\c, -0.8*\r) {\off};
    }}}

%%==========================================================================
%% 4. PARSE TREES   (M01)
%%==========================================================================
%% Use with the standard TikZ tree syntax:
%%   \begin{tikzpicture}[csftree]
%%     \node{E} child {node{E}} child {node{+}} child {node{T}};
%%   \end{tikzpicture}
\tikzset{
  csftree/.style = {
    level distance      = 9mm,
    sibling distance    = 16mm,
    every node/.style   = {font=\small},
    edge from parent/.style = {draw=bitsblue!70, thin,
                               edge from parent path=
                               {(\tikzparentnode) -- (\tikzchildnode)}},
    level 1/.style = {sibling distance=26mm},
    level 2/.style = {sibling distance=15mm},
    level 3/.style = {sibling distance=11mm},
  },
  nonterm/.style = {font=\small\itshape, color=bitsblue},
  termnode/.style = {font=\small\ttfamily, color=black},
}

%%==========================================================================
%% 5. VTABLES AND DYNAMIC DISPATCH   (M14)
%%==========================================================================
%% \vtable{<name>}{<x>}{<y>}{<class label>}{<comma-list of method slots>}
\newcommand{\vtable}[5]{%
  \begin{scope}[shift={(#2,#3)}]
    \node[csfhead, anchor=north west, minimum width=30mm] (#1head)
      at (0,0) {\scriptsize #4};
    \foreach \m [count=\i from 0] in {#5} {
      \node[csfslot, anchor=north west, minimum width=30mm] (#1\i)
        at (0, -0.6-0.55*\i) {\m};
    }
  \end{scope}}

%% \objptr{<name>}{<x>}{<y>}{<label>} -- an object header holding a vptr
\newcommand{\objptr}[4]{%
  \node[csfbox, anchor=north west, minimum width=20mm] (#1) at (#2,#3) {#4};}

%%==========================================================================
%% 6. REDUCTION SEQUENCES   (M12, M13)
%%==========================================================================
%% \redstep{<term>}{<annotation>} -- one line of a beta-reduction trace
\newcommand{\redstep}[2]{%
  \noindent$#1$ \hfill {\csflabelfont #2}\par}
\newcommand{\csflabelfont}{\scriptsize\sffamily\color{mutedgrey}}

%% Confluence diamond for Church-Rosser
%% \diamondCR{<top>}{<left>}{<right>}{<bottom>}
\newcommand{\diamondCR}[4]{%
  \node (t) at (0,2)      {$#1$};
  \node (l) at (-1.8,1)   {$#2$};
  \node (r) at (1.8,1)    {$#3$};
  \node (b) at (0,0)      {$#4$};
  \draw[csfarrow] (t) -- (l);
  \draw[csfarrow] (t) -- (r);
  \draw[csfarrow, dashed] (l) -- (b);
  \draw[csfarrow, dashed] (r) -- (b);}

%%==========================================================================
%% 7. UNIFICATION AND SEARCH TREES   (M15)
%%==========================================================================
\tikzset{
  csfsld/.style = {
    level distance = 12mm,
    sibling distance = 24mm,
    every node/.style = {font=\scriptsize\ttfamily},
    edge from parent/.style = {draw=bitsblue!70, thin,
      edge from parent path={(\tikzparentnode) -- (\tikzchildnode)}},
  },
  csffail/.style = {font=\scriptsize\sffamily, color=errred},
  csfsucc/.style = {font=\scriptsize\sffamily, color=okgreen},
}

%%==========================================================================
%% 8. SCOPE AND BINDING CHAINS   (M02, M09)
%%==========================================================================
%% \scopebox{<name>}{<x>}{<y>}{<w>}{<h>}{<label>} -- a nested scope rectangle
\newcommand{\scopebox}[6]{%
  \draw[draw=bitsblue!60, thick, rounded corners=2pt]
    (#2,#3) rectangle ++(#4,#5);
  \node[csflabel, anchor=north west] at (#2+0.08,#3+#5-0.05) {#6};}

%% \staticlink{<from>}{<to>} and \dynamiclink{<from>}{<to>}
\newcommand{\staticlink}[2]{%
  \draw[csfarrow] (#1) to[bend left=20] node[csflabel,above] {static} (#2);}
\newcommand{\dynamiclink}[2]{%
  \draw[-{Stealth[length=2mm]}, mutedgrey, thick] (#1)
    to[bend right=20] node[csflabel,below] {dynamic} (#2);}

%%==========================================================================
%% 9. GENERIC COMPARISON STRIP
%% A horizontal band of language names used to head cross-language figures.
%%==========================================================================
%% Boxes size themselves to their text and chain left to right, so long
%% labels are never clipped.
\newcommand{\langstrip}[1]{%
  \coordinate (strip@prev) at (0,0);
  \foreach \l [count=\i from 0] in {#1} {
    \ifnum\i=0
      \node[draw=bitsblue!40, fill=bitsblue!6, minimum height=8mm,
            inner xsep=3mm, font=\small\sffamily, anchor=west]
        (strip@prev) at (0,0) {\l};
    \else
      \node[draw=bitsblue!40, fill=bitsblue!6, minimum height=8mm,
            inner xsep=3mm, font=\small\sffamily, anchor=west,
            right=2mm of strip@prev] (strip@new) {\l};
      \coordinate (strip@prev) at (strip@new.east);
    \fi}}
 once, from the preamble of slides or notes.
%%
%% Every macro here produces a figure a student can redraw by hand on an
%% answer sheet in under two minutes. If a new macro cannot meet that test,
%% it does not belong in this file.
%%%%%%%%%%%%%%%%%%%%%%%%%%%%%%%%%%%%%%%%%%%%%%%%%%%%%%%%%%%%%%%%%%%%%%%%%%%%%%%

\tikzset{
  csfbox/.style   = {draw=bitsblue, thick, fill=white,
                     minimum height=6mm, minimum width=22mm, font=\small\ttfamily},
  csfslot/.style  = {draw=bitsblue!70, fill=white,
                     minimum height=5.5mm, minimum width=16mm, font=\scriptsize\ttfamily},
  csfhead/.style  = {draw=bitsblue, fill=bitsblue!12, thick,
                     minimum height=6mm, minimum width=22mm,
                     font=\small\sffamily\bfseries},
  csfdead/.style  = {draw=errred, fill=errred!8, thick,
                     minimum height=6mm, minimum width=22mm, font=\small\ttfamily},
  csflive/.style  = {draw=okgreen, fill=okgreen!8, thick,
                     minimum height=6mm, minimum width=22mm, font=\small\ttfamily},
  csflabel/.style = {font=\scriptsize\sffamily, color=mutedgrey},
  csfarrow/.style = {-{Stealth[length=2mm]}, bitsblue, thick},
  csfdash/.style  = {-{Stealth[length=2mm]}, mutedgrey, dashed, thin},
  csfnode/.style  = {draw=bitsblue, circle, thick, fill=white,
                     inner sep=1.6pt, font=\small},
  csfleaf/.style  = {draw=none, font=\small\ttfamily, color=black},
}

%%==========================================================================
%% 1. ACTIVATION RECORDS AND THE CALL STACK   (M09, M10)
%%==========================================================================
%% \arframe{<x>}{<y>}{<name>}{<comma-list of slot labels, bottom to top>}
%% Draws one activation record with a titled header and stacked slots.
\newcounter{csf@slotcount}
\newcommand{\arframe}[4]{%
  \setcounter{csf@slotcount}{0}%
  \foreach \slot in {#4} {\stepcounter{csf@slotcount}}%
  \begin{scope}[shift={(#1,#2)}]
    \foreach \slot [count=\i from 0] in {#4} {
      \node[csfslot, anchor=south west] at (0, 0.55*\i) {\slot};
    }
    \node[csfhead, anchor=south west, minimum width=16mm]
      at (0, {0.55*\value{csf@slotcount}}) {\scriptsize #3};
  \end{scope}}

%% \stacklabel{<x>}{<y>}{<text>} -- annotate a frame from the side
\newcommand{\stacklabel}[3]{\node[csflabel, anchor=west] at (#1,#2) {#3};}

%% \stackgrows{<x>}{<ytop>}{<ybot>} -- the downward growth arrow
\newcommand{\stackgrows}[3]{%
  \draw[csfdash] (#1,#2) -- (#1,#3)
    node[midway, csflabel, rotate=90, anchor=south, yshift=2pt] {stack grows};}

%%==========================================================================
%% 2. HEAP BLOCKS AND POINTERS   (M07)
%%==========================================================================
%% \heapblock{<name>}{<x>}{<y>}{<content>}{<state: live|dead|plain>}
\newcommand{\heapblock}[5]{%
  \ifthenelse{\equal{#5}{dead}}%
    {\node[csfdead, anchor=south west] (#1) at (#2,#3) {#4};}%
    {\ifthenelse{\equal{#5}{live}}%
      {\node[csflive, anchor=south west] (#1) at (#2,#3) {#4};}%
      {\node[csfbox, anchor=south west] (#1) at (#2,#3) {#4};}}}

%% \pointsto{<from>}{<to>} and \dangling{<from>}
\newcommand{\pointsto}[2]{\draw[csfarrow] (#1) -- (#2);}
\newcommand{\dangling}[1]{%
  \draw[-{Stealth[length=2mm]}, errred, thick]
    (#1) -- ++(1.2,0) node[anchor=west, csflabel, color=errred] {dangling};}
\newcommand{\unreachable}[1]{%
  \node[csflabel, color=errred, anchor=west, xshift=3mm] at (#1.east)
    {unreachable};}

%%==========================================================================
%% 3. ARRAY LAYOUT AND ADDRESS COMPUTATION   (M05)
%%==========================================================================
%% \memrow{<x>}{<y>}{<n cells>}{<first index>}{<colour>}
%% A run of contiguous memory cells with index labels beneath.
\newcommand{\memrow}[5]{%
  \begin{scope}[shift={(#1,#2)}]
    \foreach \i in {0,...,\numexpr#3-1\relax} {
      \node[draw=#5, fill=#5!10, minimum size=7mm, font=\scriptsize\ttfamily,
            anchor=south west] at (0.7*\i, 0) {};
      \node[csflabel, anchor=north] at (0.7*\i+0.35, -0.05)
        {\pgfmathparse{int(\i+#4)}\pgfmathresult};
    }
  \end{scope}}

%% \arraycell{<x>}{<y>}{<label>}{<colour>} -- one highlighted cell
\newcommand{\arraycell}[4]{%
  \node[draw=#4, fill=#4!25, thick, minimum size=7mm,
        font=\scriptsize\ttfamily, anchor=south west] at (#1,#2) {#3};}

%% \rowmajorgrid{<rows>}{<cols>} -- grid annotated with linear offsets
\newcommand{\rowmajorgrid}[2]{%
  \foreach \r in {0,...,\numexpr#1-1\relax} {
    \foreach \c in {0,...,\numexpr#2-1\relax} {
      \pgfmathtruncatemacro{\off}{\r*#2+\c}
      \node[draw=bitsblue!70, fill=bitsblue!6, minimum size=8mm,
            font=\scriptsize\ttfamily, anchor=north west]
        at (0.8*\c, -0.8*\r) {\off};
    }}}

%% \colmajorgrid{<rows>}{<cols>}
\newcommand{\colmajorgrid}[2]{%
  \foreach \r in {0,...,\numexpr#1-1\relax} {
    \foreach \c in {0,...,\numexpr#2-1\relax} {
      \pgfmathtruncatemacro{\off}{\c*#1+\r}
      \node[draw=bitsgold!80!black, fill=bitsgold!10, minimum size=8mm,
            font=\scriptsize\ttfamily, anchor=north west]
        at (0.8*\c, -0.8*\r) {\off};
    }}}

%%==========================================================================
%% 4. PARSE TREES   (M01)
%%==========================================================================
%% Use with the standard TikZ tree syntax:
%%   \begin{tikzpicture}[csftree]
%%     \node{E} child {node{E}} child {node{+}} child {node{T}};
%%   \end{tikzpicture}
\tikzset{
  csftree/.style = {
    level distance      = 9mm,
    sibling distance    = 16mm,
    every node/.style   = {font=\small},
    edge from parent/.style = {draw=bitsblue!70, thin,
                               edge from parent path=
                               {(\tikzparentnode) -- (\tikzchildnode)}},
    level 1/.style = {sibling distance=26mm},
    level 2/.style = {sibling distance=15mm},
    level 3/.style = {sibling distance=11mm},
  },
  nonterm/.style = {font=\small\itshape, color=bitsblue},
  termnode/.style = {font=\small\ttfamily, color=black},
}

%%==========================================================================
%% 5. VTABLES AND DYNAMIC DISPATCH   (M14)
%%==========================================================================
%% \vtable{<name>}{<x>}{<y>}{<class label>}{<comma-list of method slots>}
\newcommand{\vtable}[5]{%
  \begin{scope}[shift={(#2,#3)}]
    \node[csfhead, anchor=north west, minimum width=30mm] (#1head)
      at (0,0) {\scriptsize #4};
    \foreach \m [count=\i from 0] in {#5} {
      \node[csfslot, anchor=north west, minimum width=30mm] (#1\i)
        at (0, -0.6-0.55*\i) {\m};
    }
  \end{scope}}

%% \objptr{<name>}{<x>}{<y>}{<label>} -- an object header holding a vptr
\newcommand{\objptr}[4]{%
  \node[csfbox, anchor=north west, minimum width=20mm] (#1) at (#2,#3) {#4};}

%%==========================================================================
%% 6. REDUCTION SEQUENCES   (M12, M13)
%%==========================================================================
%% \redstep{<term>}{<annotation>} -- one line of a beta-reduction trace
\newcommand{\redstep}[2]{%
  \noindent$#1$ \hfill {\csflabelfont #2}\par}
\newcommand{\csflabelfont}{\scriptsize\sffamily\color{mutedgrey}}

%% Confluence diamond for Church-Rosser
%% \diamondCR{<top>}{<left>}{<right>}{<bottom>}
\newcommand{\diamondCR}[4]{%
  \node (t) at (0,2)      {$#1$};
  \node (l) at (-1.8,1)   {$#2$};
  \node (r) at (1.8,1)    {$#3$};
  \node (b) at (0,0)      {$#4$};
  \draw[csfarrow] (t) -- (l);
  \draw[csfarrow] (t) -- (r);
  \draw[csfarrow, dashed] (l) -- (b);
  \draw[csfarrow, dashed] (r) -- (b);}

%%==========================================================================
%% 7. UNIFICATION AND SEARCH TREES   (M15)
%%==========================================================================
\tikzset{
  csfsld/.style = {
    level distance = 12mm,
    sibling distance = 24mm,
    every node/.style = {font=\scriptsize\ttfamily},
    edge from parent/.style = {draw=bitsblue!70, thin,
      edge from parent path={(\tikzparentnode) -- (\tikzchildnode)}},
  },
  csffail/.style = {font=\scriptsize\sffamily, color=errred},
  csfsucc/.style = {font=\scriptsize\sffamily, color=okgreen},
}

%%==========================================================================
%% 8. SCOPE AND BINDING CHAINS   (M02, M09)
%%==========================================================================
%% \scopebox{<name>}{<x>}{<y>}{<w>}{<h>}{<label>} -- a nested scope rectangle
\newcommand{\scopebox}[6]{%
  \draw[draw=bitsblue!60, thick, rounded corners=2pt]
    (#2,#3) rectangle ++(#4,#5);
  \node[csflabel, anchor=north west] at (#2+0.08,#3+#5-0.05) {#6};}

%% \staticlink{<from>}{<to>} and \dynamiclink{<from>}{<to>}
\newcommand{\staticlink}[2]{%
  \draw[csfarrow] (#1) to[bend left=20] node[csflabel,above] {static} (#2);}
\newcommand{\dynamiclink}[2]{%
  \draw[-{Stealth[length=2mm]}, mutedgrey, thick] (#1)
    to[bend right=20] node[csflabel,below] {dynamic} (#2);}

%%==========================================================================
%% 9. GENERIC COMPARISON STRIP
%% A horizontal band of language names used to head cross-language figures.
%%==========================================================================
%% Boxes size themselves to their text and chain left to right, so long
%% labels are never clipped.
\newcommand{\langstrip}[1]{%
  \coordinate (strip@prev) at (0,0);
  \foreach \l [count=\i from 0] in {#1} {
    \ifnum\i=0
      \node[draw=bitsblue!40, fill=bitsblue!6, minimum height=8mm,
            inner xsep=3mm, font=\small\sffamily, anchor=west]
        (strip@prev) at (0,0) {\l};
    \else
      \node[draw=bitsblue!40, fill=bitsblue!6, minimum height=8mm,
            inner xsep=3mm, font=\small\sffamily, anchor=west,
            right=2mm of strip@prev] (strip@new) {\l};
      \coordinate (strip@prev) at (strip@new.east);
    \fi}}
 once, from the preamble of slides or notes.
%%
%% Every macro here produces a figure a student can redraw by hand on an
%% answer sheet in under two minutes. If a new macro cannot meet that test,
%% it does not belong in this file.
%%%%%%%%%%%%%%%%%%%%%%%%%%%%%%%%%%%%%%%%%%%%%%%%%%%%%%%%%%%%%%%%%%%%%%%%%%%%%%%

\tikzset{
  csfbox/.style   = {draw=bitsblue, thick, fill=white,
                     minimum height=6mm, minimum width=22mm, font=\small\ttfamily},
  csfslot/.style  = {draw=bitsblue!70, fill=white,
                     minimum height=5.5mm, minimum width=16mm, font=\scriptsize\ttfamily},
  csfhead/.style  = {draw=bitsblue, fill=bitsblue!12, thick,
                     minimum height=6mm, minimum width=22mm,
                     font=\small\sffamily\bfseries},
  csfdead/.style  = {draw=errred, fill=errred!8, thick,
                     minimum height=6mm, minimum width=22mm, font=\small\ttfamily},
  csflive/.style  = {draw=okgreen, fill=okgreen!8, thick,
                     minimum height=6mm, minimum width=22mm, font=\small\ttfamily},
  csflabel/.style = {font=\scriptsize\sffamily, color=mutedgrey},
  csfarrow/.style = {-{Stealth[length=2mm]}, bitsblue, thick},
  csfdash/.style  = {-{Stealth[length=2mm]}, mutedgrey, dashed, thin},
  csfnode/.style  = {draw=bitsblue, circle, thick, fill=white,
                     inner sep=1.6pt, font=\small},
  csfleaf/.style  = {draw=none, font=\small\ttfamily, color=black},
}

%%==========================================================================
%% 1. ACTIVATION RECORDS AND THE CALL STACK   (M09, M10)
%%==========================================================================
%% \arframe{<x>}{<y>}{<name>}{<comma-list of slot labels, bottom to top>}
%% Draws one activation record with a titled header and stacked slots.
\newcounter{csf@slotcount}
\newcommand{\arframe}[4]{%
  \setcounter{csf@slotcount}{0}%
  \foreach \slot in {#4} {\stepcounter{csf@slotcount}}%
  \begin{scope}[shift={(#1,#2)}]
    \foreach \slot [count=\i from 0] in {#4} {
      \node[csfslot, anchor=south west] at (0, 0.55*\i) {\slot};
    }
    \node[csfhead, anchor=south west, minimum width=16mm]
      at (0, {0.55*\value{csf@slotcount}}) {\scriptsize #3};
  \end{scope}}

%% \stacklabel{<x>}{<y>}{<text>} -- annotate a frame from the side
\newcommand{\stacklabel}[3]{\node[csflabel, anchor=west] at (#1,#2) {#3};}

%% \stackgrows{<x>}{<ytop>}{<ybot>} -- the downward growth arrow
\newcommand{\stackgrows}[3]{%
  \draw[csfdash] (#1,#2) -- (#1,#3)
    node[midway, csflabel, rotate=90, anchor=south, yshift=2pt] {stack grows};}

%%==========================================================================
%% 2. HEAP BLOCKS AND POINTERS   (M07)
%%==========================================================================
%% \heapblock{<name>}{<x>}{<y>}{<content>}{<state: live|dead|plain>}
\newcommand{\heapblock}[5]{%
  \ifthenelse{\equal{#5}{dead}}%
    {\node[csfdead, anchor=south west] (#1) at (#2,#3) {#4};}%
    {\ifthenelse{\equal{#5}{live}}%
      {\node[csflive, anchor=south west] (#1) at (#2,#3) {#4};}%
      {\node[csfbox, anchor=south west] (#1) at (#2,#3) {#4};}}}

%% \pointsto{<from>}{<to>} and \dangling{<from>}
\newcommand{\pointsto}[2]{\draw[csfarrow] (#1) -- (#2);}
\newcommand{\dangling}[1]{%
  \draw[-{Stealth[length=2mm]}, errred, thick]
    (#1) -- ++(1.2,0) node[anchor=west, csflabel, color=errred] {dangling};}
\newcommand{\unreachable}[1]{%
  \node[csflabel, color=errred, anchor=west, xshift=3mm] at (#1.east)
    {unreachable};}

%%==========================================================================
%% 3. ARRAY LAYOUT AND ADDRESS COMPUTATION   (M05)
%%==========================================================================
%% \memrow{<x>}{<y>}{<n cells>}{<first index>}{<colour>}
%% A run of contiguous memory cells with index labels beneath.
\newcommand{\memrow}[5]{%
  \begin{scope}[shift={(#1,#2)}]
    \foreach \i in {0,...,\numexpr#3-1\relax} {
      \node[draw=#5, fill=#5!10, minimum size=7mm, font=\scriptsize\ttfamily,
            anchor=south west] at (0.7*\i, 0) {};
      \node[csflabel, anchor=north] at (0.7*\i+0.35, -0.05)
        {\pgfmathparse{int(\i+#4)}\pgfmathresult};
    }
  \end{scope}}

%% \arraycell{<x>}{<y>}{<label>}{<colour>} -- one highlighted cell
\newcommand{\arraycell}[4]{%
  \node[draw=#4, fill=#4!25, thick, minimum size=7mm,
        font=\scriptsize\ttfamily, anchor=south west] at (#1,#2) {#3};}

%% \rowmajorgrid{<rows>}{<cols>} -- grid annotated with linear offsets
\newcommand{\rowmajorgrid}[2]{%
  \foreach \r in {0,...,\numexpr#1-1\relax} {
    \foreach \c in {0,...,\numexpr#2-1\relax} {
      \pgfmathtruncatemacro{\off}{\r*#2+\c}
      \node[draw=bitsblue!70, fill=bitsblue!6, minimum size=8mm,
            font=\scriptsize\ttfamily, anchor=north west]
        at (0.8*\c, -0.8*\r) {\off};
    }}}

%% \colmajorgrid{<rows>}{<cols>}
\newcommand{\colmajorgrid}[2]{%
  \foreach \r in {0,...,\numexpr#1-1\relax} {
    \foreach \c in {0,...,\numexpr#2-1\relax} {
      \pgfmathtruncatemacro{\off}{\c*#1+\r}
      \node[draw=bitsgold!80!black, fill=bitsgold!10, minimum size=8mm,
            font=\scriptsize\ttfamily, anchor=north west]
        at (0.8*\c, -0.8*\r) {\off};
    }}}

%%==========================================================================
%% 4. PARSE TREES   (M01)
%%==========================================================================
%% Use with the standard TikZ tree syntax:
%%   \begin{tikzpicture}[csftree]
%%     \node{E} child {node{E}} child {node{+}} child {node{T}};
%%   \end{tikzpicture}
\tikzset{
  csftree/.style = {
    level distance      = 9mm,
    sibling distance    = 16mm,
    every node/.style   = {font=\small},
    edge from parent/.style = {draw=bitsblue!70, thin,
                               edge from parent path=
                               {(\tikzparentnode) -- (\tikzchildnode)}},
    level 1/.style = {sibling distance=26mm},
    level 2/.style = {sibling distance=15mm},
    level 3/.style = {sibling distance=11mm},
  },
  nonterm/.style = {font=\small\itshape, color=bitsblue},
  termnode/.style = {font=\small\ttfamily, color=black},
}

%%==========================================================================
%% 5. VTABLES AND DYNAMIC DISPATCH   (M14)
%%==========================================================================
%% \vtable{<name>}{<x>}{<y>}{<class label>}{<comma-list of method slots>}
\newcommand{\vtable}[5]{%
  \begin{scope}[shift={(#2,#3)}]
    \node[csfhead, anchor=north west, minimum width=30mm] (#1head)
      at (0,0) {\scriptsize #4};
    \foreach \m [count=\i from 0] in {#5} {
      \node[csfslot, anchor=north west, minimum width=30mm] (#1\i)
        at (0, -0.6-0.55*\i) {\m};
    }
  \end{scope}}

%% \objptr{<name>}{<x>}{<y>}{<label>} -- an object header holding a vptr
\newcommand{\objptr}[4]{%
  \node[csfbox, anchor=north west, minimum width=20mm] (#1) at (#2,#3) {#4};}

%%==========================================================================
%% 6. REDUCTION SEQUENCES   (M12, M13)
%%==========================================================================
%% \redstep{<term>}{<annotation>} -- one line of a beta-reduction trace
\newcommand{\redstep}[2]{%
  \noindent$#1$ \hfill {\csflabelfont #2}\par}
\newcommand{\csflabelfont}{\scriptsize\sffamily\color{mutedgrey}}

%% Confluence diamond for Church-Rosser
%% \diamondCR{<top>}{<left>}{<right>}{<bottom>}
\newcommand{\diamondCR}[4]{%
  \node (t) at (0,2)      {$#1$};
  \node (l) at (-1.8,1)   {$#2$};
  \node (r) at (1.8,1)    {$#3$};
  \node (b) at (0,0)      {$#4$};
  \draw[csfarrow] (t) -- (l);
  \draw[csfarrow] (t) -- (r);
  \draw[csfarrow, dashed] (l) -- (b);
  \draw[csfarrow, dashed] (r) -- (b);}

%%==========================================================================
%% 7. UNIFICATION AND SEARCH TREES   (M15)
%%==========================================================================
\tikzset{
  csfsld/.style = {
    level distance = 12mm,
    sibling distance = 24mm,
    every node/.style = {font=\scriptsize\ttfamily},
    edge from parent/.style = {draw=bitsblue!70, thin,
      edge from parent path={(\tikzparentnode) -- (\tikzchildnode)}},
  },
  csffail/.style = {font=\scriptsize\sffamily, color=errred},
  csfsucc/.style = {font=\scriptsize\sffamily, color=okgreen},
}

%%==========================================================================
%% 8. SCOPE AND BINDING CHAINS   (M02, M09)
%%==========================================================================
%% \scopebox{<name>}{<x>}{<y>}{<w>}{<h>}{<label>} -- a nested scope rectangle
\newcommand{\scopebox}[6]{%
  \draw[draw=bitsblue!60, thick, rounded corners=2pt]
    (#2,#3) rectangle ++(#4,#5);
  \node[csflabel, anchor=north west] at (#2+0.08,#3+#5-0.05) {#6};}

%% \staticlink{<from>}{<to>} and \dynamiclink{<from>}{<to>}
\newcommand{\staticlink}[2]{%
  \draw[csfarrow] (#1) to[bend left=20] node[csflabel,above] {static} (#2);}
\newcommand{\dynamiclink}[2]{%
  \draw[-{Stealth[length=2mm]}, mutedgrey, thick] (#1)
    to[bend right=20] node[csflabel,below] {dynamic} (#2);}

%%==========================================================================
%% 9. GENERIC COMPARISON STRIP
%% A horizontal band of language names used to head cross-language figures.
%%==========================================================================
%% Boxes size themselves to their text and chain left to right, so long
%% labels are never clipped.
\newcommand{\langstrip}[1]{%
  \coordinate (strip@prev) at (0,0);
  \foreach \l [count=\i from 0] in {#1} {
    \ifnum\i=0
      \node[draw=bitsblue!40, fill=bitsblue!6, minimum height=8mm,
            inner xsep=3mm, font=\small\sffamily, anchor=west]
        (strip@prev) at (0,0) {\l};
    \else
      \node[draw=bitsblue!40, fill=bitsblue!6, minimum height=8mm,
            inner xsep=3mm, font=\small\sffamily, anchor=west,
            right=2mm of strip@prev] (strip@new) {\l};
      \coordinate (strip@prev) at (strip@new.east);
    \fi}}
 once, from the preamble of slides or notes.
%%
%% Every macro here produces a figure a student can redraw by hand on an
%% answer sheet in under two minutes. If a new macro cannot meet that test,
%% it does not belong in this file.
%%%%%%%%%%%%%%%%%%%%%%%%%%%%%%%%%%%%%%%%%%%%%%%%%%%%%%%%%%%%%%%%%%%%%%%%%%%%%%%

\tikzset{
  csfbox/.style   = {draw=bitsblue, thick, fill=white,
                     minimum height=6mm, minimum width=22mm, font=\small\ttfamily},
  csfslot/.style  = {draw=bitsblue!70, fill=white,
                     minimum height=5.5mm, minimum width=16mm, font=\scriptsize\ttfamily},
  csfhead/.style  = {draw=bitsblue, fill=bitsblue!12, thick,
                     minimum height=6mm, minimum width=22mm,
                     font=\small\sffamily\bfseries},
  csfdead/.style  = {draw=errred, fill=errred!8, thick,
                     minimum height=6mm, minimum width=22mm, font=\small\ttfamily},
  csflive/.style  = {draw=okgreen, fill=okgreen!8, thick,
                     minimum height=6mm, minimum width=22mm, font=\small\ttfamily},
  csflabel/.style = {font=\scriptsize\sffamily, color=mutedgrey},
  csfarrow/.style = {-{Stealth[length=2mm]}, bitsblue, thick},
  csfdash/.style  = {-{Stealth[length=2mm]}, mutedgrey, dashed, thin},
  csfnode/.style  = {draw=bitsblue, circle, thick, fill=white,
                     inner sep=1.6pt, font=\small},
  csfleaf/.style  = {draw=none, font=\small\ttfamily, color=black},
}

%%==========================================================================
%% 1. ACTIVATION RECORDS AND THE CALL STACK   (M09, M10)
%%==========================================================================
%% \arframe{<x>}{<y>}{<name>}{<comma-list of slot labels, bottom to top>}
%% Draws one activation record with a titled header and stacked slots.
\newcounter{csf@slotcount}
\newcommand{\arframe}[4]{%
  \setcounter{csf@slotcount}{0}%
  \foreach \slot in {#4} {\stepcounter{csf@slotcount}}%
  \begin{scope}[shift={(#1,#2)}]
    \foreach \slot [count=\i from 0] in {#4} {
      \node[csfslot, anchor=south west] at (0, 0.55*\i) {\slot};
    }
    \node[csfhead, anchor=south west, minimum width=16mm]
      at (0, {0.55*\value{csf@slotcount}}) {\scriptsize #3};
  \end{scope}}

%% \stacklabel{<x>}{<y>}{<text>} -- annotate a frame from the side
\newcommand{\stacklabel}[3]{\node[csflabel, anchor=west] at (#1,#2) {#3};}

%% \stackgrows{<x>}{<ytop>}{<ybot>} -- the downward growth arrow
\newcommand{\stackgrows}[3]{%
  \draw[csfdash] (#1,#2) -- (#1,#3)
    node[midway, csflabel, rotate=90, anchor=south, yshift=2pt] {stack grows};}

%%==========================================================================
%% 2. HEAP BLOCKS AND POINTERS   (M07)
%%==========================================================================
%% \heapblock{<name>}{<x>}{<y>}{<content>}{<state: live|dead|plain>}
\newcommand{\heapblock}[5]{%
  \ifthenelse{\equal{#5}{dead}}%
    {\node[csfdead, anchor=south west] (#1) at (#2,#3) {#4};}%
    {\ifthenelse{\equal{#5}{live}}%
      {\node[csflive, anchor=south west] (#1) at (#2,#3) {#4};}%
      {\node[csfbox, anchor=south west] (#1) at (#2,#3) {#4};}}}

%% \pointsto{<from>}{<to>} and \dangling{<from>}
\newcommand{\pointsto}[2]{\draw[csfarrow] (#1) -- (#2);}
\newcommand{\dangling}[1]{%
  \draw[-{Stealth[length=2mm]}, errred, thick]
    (#1) -- ++(1.2,0) node[anchor=west, csflabel, color=errred] {dangling};}
\newcommand{\unreachable}[1]{%
  \node[csflabel, color=errred, anchor=west, xshift=3mm] at (#1.east)
    {unreachable};}

%%==========================================================================
%% 3. ARRAY LAYOUT AND ADDRESS COMPUTATION   (M05)
%%==========================================================================
%% \memrow{<x>}{<y>}{<n cells>}{<first index>}{<colour>}
%% A run of contiguous memory cells with index labels beneath.
\newcommand{\memrow}[5]{%
  \begin{scope}[shift={(#1,#2)}]
    \foreach \i in {0,...,\numexpr#3-1\relax} {
      \node[draw=#5, fill=#5!10, minimum size=7mm, font=\scriptsize\ttfamily,
            anchor=south west] at (0.7*\i, 0) {};
      \node[csflabel, anchor=north] at (0.7*\i+0.35, -0.05)
        {\pgfmathparse{int(\i+#4)}\pgfmathresult};
    }
  \end{scope}}

%% \arraycell{<x>}{<y>}{<label>}{<colour>} -- one highlighted cell
\newcommand{\arraycell}[4]{%
  \node[draw=#4, fill=#4!25, thick, minimum size=7mm,
        font=\scriptsize\ttfamily, anchor=south west] at (#1,#2) {#3};}

%% \rowmajorgrid{<rows>}{<cols>} -- grid annotated with linear offsets
\newcommand{\rowmajorgrid}[2]{%
  \foreach \r in {0,...,\numexpr#1-1\relax} {
    \foreach \c in {0,...,\numexpr#2-1\relax} {
      \pgfmathtruncatemacro{\off}{\r*#2+\c}
      \node[draw=bitsblue!70, fill=bitsblue!6, minimum size=8mm,
            font=\scriptsize\ttfamily, anchor=north west]
        at (0.8*\c, -0.8*\r) {\off};
    }}}

%% \colmajorgrid{<rows>}{<cols>}
\newcommand{\colmajorgrid}[2]{%
  \foreach \r in {0,...,\numexpr#1-1\relax} {
    \foreach \c in {0,...,\numexpr#2-1\relax} {
      \pgfmathtruncatemacro{\off}{\c*#1+\r}
      \node[draw=bitsgold!80!black, fill=bitsgold!10, minimum size=8mm,
            font=\scriptsize\ttfamily, anchor=north west]
        at (0.8*\c, -0.8*\r) {\off};
    }}}

%%==========================================================================
%% 4. PARSE TREES   (M01)
%%==========================================================================
%% Use with the standard TikZ tree syntax:
%%   \begin{tikzpicture}[csftree]
%%     \node{E} child {node{E}} child {node{+}} child {node{T}};
%%   \end{tikzpicture}
\tikzset{
  csftree/.style = {
    level distance      = 9mm,
    sibling distance    = 16mm,
    every node/.style   = {font=\small},
    edge from parent/.style = {draw=bitsblue!70, thin,
                               edge from parent path=
                               {(\tikzparentnode) -- (\tikzchildnode)}},
    level 1/.style = {sibling distance=26mm},
    level 2/.style = {sibling distance=15mm},
    level 3/.style = {sibling distance=11mm},
  },
  nonterm/.style = {font=\small\itshape, color=bitsblue},
  termnode/.style = {font=\small\ttfamily, color=black},
}

%%==========================================================================
%% 5. VTABLES AND DYNAMIC DISPATCH   (M14)
%%==========================================================================
%% \vtable{<name>}{<x>}{<y>}{<class label>}{<comma-list of method slots>}
\newcommand{\vtable}[5]{%
  \begin{scope}[shift={(#2,#3)}]
    \node[csfhead, anchor=north west, minimum width=30mm] (#1head)
      at (0,0) {\scriptsize #4};
    \foreach \m [count=\i from 0] in {#5} {
      \node[csfslot, anchor=north west, minimum width=30mm] (#1\i)
        at (0, -0.6-0.55*\i) {\m};
    }
  \end{scope}}

%% \objptr{<name>}{<x>}{<y>}{<label>} -- an object header holding a vptr
\newcommand{\objptr}[4]{%
  \node[csfbox, anchor=north west, minimum width=20mm] (#1) at (#2,#3) {#4};}

%%==========================================================================
%% 6. REDUCTION SEQUENCES   (M12, M13)
%%==========================================================================
%% \redstep{<term>}{<annotation>} -- one line of a beta-reduction trace
\newcommand{\redstep}[2]{%
  \noindent$#1$ \hfill {\csflabelfont #2}\par}
\newcommand{\csflabelfont}{\scriptsize\sffamily\color{mutedgrey}}

%% Confluence diamond for Church-Rosser
%% \diamondCR{<top>}{<left>}{<right>}{<bottom>}
\newcommand{\diamondCR}[4]{%
  \node (t) at (0,2)      {$#1$};
  \node (l) at (-1.8,1)   {$#2$};
  \node (r) at (1.8,1)    {$#3$};
  \node (b) at (0,0)      {$#4$};
  \draw[csfarrow] (t) -- (l);
  \draw[csfarrow] (t) -- (r);
  \draw[csfarrow, dashed] (l) -- (b);
  \draw[csfarrow, dashed] (r) -- (b);}

%%==========================================================================
%% 7. UNIFICATION AND SEARCH TREES   (M15)
%%==========================================================================
\tikzset{
  csfsld/.style = {
    level distance = 12mm,
    sibling distance = 24mm,
    every node/.style = {font=\scriptsize\ttfamily},
    edge from parent/.style = {draw=bitsblue!70, thin,
      edge from parent path={(\tikzparentnode) -- (\tikzchildnode)}},
  },
  csffail/.style = {font=\scriptsize\sffamily, color=errred},
  csfsucc/.style = {font=\scriptsize\sffamily, color=okgreen},
}

%%==========================================================================
%% 8. SCOPE AND BINDING CHAINS   (M02, M09)
%%==========================================================================
%% \scopebox{<name>}{<x>}{<y>}{<w>}{<h>}{<label>} -- a nested scope rectangle
\newcommand{\scopebox}[6]{%
  \draw[draw=bitsblue!60, thick, rounded corners=2pt]
    (#2,#3) rectangle ++(#4,#5);
  \node[csflabel, anchor=north west] at (#2+0.08,#3+#5-0.05) {#6};}

%% \staticlink{<from>}{<to>} and \dynamiclink{<from>}{<to>}
\newcommand{\staticlink}[2]{%
  \draw[csfarrow] (#1) to[bend left=20] node[csflabel,above] {static} (#2);}
\newcommand{\dynamiclink}[2]{%
  \draw[-{Stealth[length=2mm]}, mutedgrey, thick] (#1)
    to[bend right=20] node[csflabel,below] {dynamic} (#2);}

%%==========================================================================
%% 9. GENERIC COMPARISON STRIP
%% A horizontal band of language names used to head cross-language figures.
%%==========================================================================
%% Boxes size themselves to their text and chain left to right, so long
%% labels are never clipped.
\newcommand{\langstrip}[1]{%
  \coordinate (strip@prev) at (0,0);
  \foreach \l [count=\i from 0] in {#1} {
    \ifnum\i=0
      \node[draw=bitsblue!40, fill=bitsblue!6, minimum height=8mm,
            inner xsep=3mm, font=\small\sffamily, anchor=west]
        (strip@prev) at (0,0) {\l};
    \else
      \node[draw=bitsblue!40, fill=bitsblue!6, minimum height=8mm,
            inner xsep=3mm, font=\small\sffamily, anchor=west,
            right=2mm of strip@prev] (strip@new) {\l};
      \coordinate (strip@prev) at (strip@new.east);
    \fi}}
